\documentclass{report}
\usepackage{amsfonts, amsmath, mathabx}
\newcommand\bR{{\mathbb R}}
\newcommand\bQ{{\mathbb Q}}
\newcommand\bN{{\mathbb N}}
\newcommand\cC{{\mathcal C}}
\newcommand\cO{{\mathcal O}}
\newcommand\ve{{\varepsilon}}
%\newcommand\0{\noindent}
\newcommand\virg{\quad,\quad}

\setlength{\marginparwidth}{8em} \setlength{\marginparsep}{1em}
\setlength{\textwidth}{34.8pc} \setlength{\textheight}{51.8pc}
\setlength{\oddsidemargin}{1truecm}
\setlength{\topmargin}{-0.5truecm}
\renewcommand{\thefootnote}{\relax}
\newcommand\barra{\vrule height.5cm depth0pt width0pt}

\begin{document}
\section*{Compito per il 15/03/2024}\footnote{ I punteggi, nel discutibile stile americano, saranno A,B,C e E, dove $E$ significa che  la soluzione non è ritenuta sufficiente. Alla fine si avranno tre punti se si sono ottenuti, almeno nei due terzi dei compiti, A; 2 punti se si sono ottenuti,  almeno nei due terzi dei compiti, A o B; un punto se si sono ottenuti, almeno nei due terzi dei compiti, A, B o C.}
\begin{enumerate}
\item Dato un intervallo $I=(a,b)\subset \bR$ si consideri la funzione $d_I:I^2\to\bR_+$ definita da
\[
d_I(x,y)=\left|\ln\frac{ (b-x)(y-a)}{(b-y)(x-a)}\right|.
\]
\begin{enumerate}
\item Si mostri che $d_I$ è una distanza.
\item Si calcoli,  $\lim_{n\to\infty} d_I(a+1/n, (a+b)/2)$.
\item Quale è il diametro di $I$ nella metrica $d_I$?
\item Data la funzione $f(x)=\frac{x-a}{b-a}$, con dominio $I$, si dica quale è l'immagine $J=f(I)$ e si calcoli il rapporto $\frac{d_J(f(x),f(y))}{d_I(x,y)}$.
\item Sia $I_n=(0,n)$, si calcoli $d_\infty(x,y)=\lim_{n\to\infty}d_{I_n}(x,y)$ e si mostri che è una metrica in $(0,\infty)$.
\item Data la funzione $f(x)=\frac{x-a}{b-x}$, con dominio $I$, si dica quale è l'immagine $J=f(I)$ e si calcoli il rapporto $\frac{d_\infty(f(x),f(y))}{d_I(x,y)}$.
\end{enumerate}
\end{enumerate}
\newpage

\centerline{\huge\bf Soluzione}
\begin{enumerate}
\item Soluzione dell'esercizio 1:
\begin{enumerate}
\item L'unica cosa che non è immediatamente ovvia è la disuguaglianza triangolare
\[
\begin{split}
&d_I(x,y)=\left|\ln\frac{ (b-x)(z-a)(b-z)(y-a)}{(x-a)(b-z)(z-a)(b-y)}\right|\\
&=\left|\ln\frac{ (b-x)(z-a)}{(x-a)(b-z)}+\frac{(b-z)(y-a)}{(z-a)(b-y)}\right|
\leq \left|\ln\frac{ (b-x)(z-a)}{(x-a)(b-z)}\right|+\left|\frac{(b-z)(y-a)}{(z-a)(b-y)}\right|\\
&=d_I(x,z)+d_I(z,y).
\end{split}
\]
\item Si noti che, se  $\frac 1n\leq b-a$, allora
\[
\begin{split}
d_I(a+1/n, (a+b)/2)&=\left|\ln\frac{ (b-a-1/n)((a+b)/2-a)}{(1/n)(b-(a+b)/2)}\right|
=\left|\ln (b-a-1/n)n\right|\\
&=\ln [(b-a)n-1].
\end{split}
\]
Quindi $\lim_{n\to\infty}d_I(a+1/n, (a+b)/2)=\infty$.
\item Visto il punto precedente il diametro è chiaramente infinito.
\item $J=(0,1)$, quindi
\[
\begin{split}
d_J(f(x),f(y))&=\left|\ln\frac{ (1-f(x))f(y)}{f(x)(1-f(y))}\right|=
\left|\ln\frac{ (1-\frac{x-a}{b-a})\frac{y-a}{b-a}}{\frac{x-a}{b-a}(1-\frac{y-a}{b-a})}\right|\\
&=\left|\ln\frac{ (b-x)(y-a)}{(x-a)(b-y)}\right|=d_I(x,y)
\end{split}
\]
ovvero, rispetto a queste metriche $f$ è una isometria.
\item Per ogni $x,y\in(0,\infty)$, per ogni $n>\max\{x,y\}$ si ha
\[
d_{I_n}(x,y)=\left|\ln\frac{ (n-x)y}{(n-y)x}\right|=\left|\ln\frac{ (1-x/n)y}{(1-y/n)x}\right|.
\]
dunque
\[
d_\infty(x,y)=\left|\ln\frac{y}{x}\right|.
\]
Si può facilmente dimostrare che è una norma argomentando come nel punto (a).
\item $J=(0,\infty)$, quindi
\[
\begin{split}
d_\infty(f(x),f(y))&=\left|\ln\frac{ f(x)}{f(y)}\right|=d_I(x,y).
\end{split}
\]
Un'altra isometria!
\end{enumerate}
\end{enumerate}
\end{document}
